\chapter{Planned Evaluation}
\label{chap:Plan}

The refined problem has not yet been run. This chapter sets out the experiments that
will test it, in the order in which they depend on one another. For each hypothesis the
decision rule is stated before the results.

\section{Pilot without a Budget}
\label{plan:sec:pilot}

The pilot isolates three of the changes of the refined problem from the fourth. It
removes the budget, raises the batch to 100 prompts per round and uses the selector of
Section~\ref{imp:sec:efficiency}, but keeps the agreement reward of the preliminary
study, because the provider term of the new objective is not yet built. It compares
uniform allocation with the selector on the three providers of
Table~\ref{imp:tab:pool}, with one seed, at the configuration of
Table~\ref{imp:tab:config}: 100 rounds of 100 prompts, a decision every 5 rounds, and
the full selection set at every decision. Each condition is run at learning rates of
$2 \times 10^{-5}$ and $10^{-4}$, which gives four runs and about 20 GPU-hours.

The pilot answers three questions. Does the student still learn with 12 times fewer
gradient steps for the same prompts, and at which learning rate? Is the gain of one
block now larger than the evaluation noise? Does the selector's learning curve separate
from uniform allocation's? At the time of writing the four runs are queued and no result
is available.

\section{Experiments}
\label{plan:sec:experiments}

\begin{description}[itemsep=4pt,topsep=4pt]
  \item[Experiment 1: the training objective.] Compare the reward term alone, the
    provider term alone and their sum (Equation~\ref{imp:eq:total}), for each provider
    separately, by learning curves and by the response-length statistics that revealed
    the collapse in the preliminary study. The objective is adopted if the sum is at
    least as good as either term alone and no run collapses.
  \item[Experiment 2: progress on MATH and code.] With the adopted objective, test
    whether the student's accuracy on MATH and code rises above its starting point by
    more than the 95\% interval of the evaluation. If it does not, repeat with a larger
    student. Without progress outside Countdown, a curriculum has only one domain worth
    choosing.
  \item[Experiment 3: non-dominance and state dependence (H1, H2).] Measure the
    marginal value of every arm by branching from shared checkpoints, as in
    Section~\ref{pre:sec:branching}, with longer branches and per-item results so that
    providers are compared on the same items within each seed. H1 is supported if, at
    some checkpoint, the best provider differs between two domains with paired intervals
    that exclude zero; H2 if, within a domain, it differs between two checkpoints.
  \item[Experiment 4: adaptive against fixed curricula (H3).] Compare the selector with
    the reference curricula of Section~\ref{imp:sec:baselines} over three seeds. H3 is
    supported if the selector reaches the final utility of the best fixed curriculum
    with fewer training prompts, with a paired interval that excludes zero.
  \item[Experiment 5: passing the teachers (H4).] Compare the student trained on
    several providers with the student trained on the best single provider, and with
    each provider's own accuracy in Table~\ref{imp:tab:pool}. H4 is supported if the
    former exceeds the best single-provider student on test utility with a paired
    interval that excludes zero.
  \item[Experiment 6: external check.] Evaluate the final students on standard
    mathematics benchmarks outside the training pool, as an out-of-distribution check.
\end{description}

\section{Schedule, Resources and Risks}
\label{plan:sec:schedule}

\begin{table}[htbp]
\centering
\small
\begin{tabular}{@{}lcccccc@{}}
\toprule
\textbf{Task} & \textbf{Nov} & \textbf{Dec} & \textbf{Jan} & \textbf{Feb} & \textbf{Mar} & \textbf{Apr} \\
\midrule
Pilot without a budget           & \textbullet &             &             &             &             &   \\
1. Training objective            & \textbullet & \textbullet &             &             &             &   \\
2. Progress on MATH and code     &             & \textbullet & \textbullet &             &             &   \\
3, 4. H1, H2 and curricula       &             &             & \textbullet & \textbullet &             &   \\
5. Passing the teachers          &             &             &             & \textbullet & \textbullet &   \\
6. External check                &             &             &             &             & \textbullet &   \\
Dissertation and presentation    &             & \textbullet &             &             & \textbullet & \textbullet \\
\bottomrule
\end{tabular}
\caption{Schedule for the remaining work.}
\label{plan:tab:schedule}
\end{table}

\paragraph{Resources.} The preliminary study used about 110 L40 GPU-hours. The remaining
work is estimated at about 210: 20 for the pilot, 40 to develop the objective, 40 for
MATH and code, 100 for the curriculum comparison over three seeds, and 10 for the
external check. These are estimates made before the pilot has run.

\paragraph{Risks.} The main risk is that learning stays confined to Countdown under the
new objective, which would leave the curriculum only one domain worth choosing and
narrow the claims. The second is that, without a budget, every reasonable curriculum
reaches the same accuracy and differs only in speed; learning curves are the headline
metric for that reason. The third is that the three providers are too alike for the
choice between them to matter, in which case a provider of another family would be
added through a text-only reward.
